\chapter{Comprendre les fonctions}
\label{exercices:08}
\section*{Exercice 8.1 -- Images et antécédents}
Soit $f(x)=x^2-4$ sur les réels. Calculer $f(-3)$, $f(0)$ et $f(2)$. Déterminer tous les antécédents de $5$, de $-4$ et de $-5$.

\section*{Exercice 8.2 -- Lire un tableau avec prudence}
Une fonction $g$ est connue par :
\[
\begin{array}{c|rrrrr}x&-2&-1&0&1&2\\\hline g(x)&3&0&-1&0&3\end{array}
\]
Donner l'image de $-1$, deux antécédents de $3$ et expliquer ce qu'on peut dire de $g(1{,}5)$.

\section*{Exercice 8.3 -- Du programme à la formule}
Un programme choisit $x$, le multiplie par $-3$, puis ajoute $7$. Écrire la fonction $h$. Chercher l'antécédent de $19$ et vérifier.

\section*{Exercice 8.4 -- Modéliser une aire}
Un rectangle a un périmètre de $30\unite{cm}$. Un côté vaut $x\unite{cm}$. Écrire son aire $A(x)$ et le domaine des valeurs possibles pour un rectangle non aplati. Calculer $A(4)$ et trouver une autre longueur donnant la même aire.
